\setcounter{section}{-1}
\section{Preliminaries}

\subsection{Basics}

For \exref{0.1.1} through \exref{0.1.4}, let
\[
    M = 
    \begin{pmatrix}
        1 & 1 \\
        0 & 1
    \end{pmatrix}
\]
and $\bb = \set{X \in \ac \mid MX = XM}$, where $\ac$ denotes the set of $2 \times 2$ matrices with real entries.

\begin{exercise}[0.1.1]
    Determine which of the following elements of $\ac$ lie in $\bb$:
    \[
        \begin{pmatrix}
            1 & 1 \\
            0 & 1
        \end{pmatrix},
        \begin{pmatrix}
            1 & 1 \\
            1 & 1
        \end{pmatrix},
        \begin{pmatrix}
            0 & 0 \\
            0 & 0
        \end{pmatrix},
        \begin{pmatrix}
            1 & 1 \\
            1 & 0
        \end{pmatrix},
        \begin{pmatrix}
            1 & 0 \\
            0 & 1
        \end{pmatrix},
        \begin{pmatrix}
            0 & 1 \\
            1 & 0 
        \end{pmatrix}
    \]
\end{exercise}

\begin{sol}
    Let $M_i$ denote the $i$-th matrix in the list.
	Then $M_1, M_3, M_5, M_6$ are $M$, the zero matrix, $I_2$, and the exchange matrix respectively.
	Among these, $M_1, M_3, M_5$ belong to $\bb$, while $M_6 \notin \bb$.
	Direct calculation shows $M_2$ and $M_4$ are not in $\bb$.
\end{sol}

\begin{exercise}[0.1.2]
    Prove that if $P, Q \in \bb$, then $P + Q \in \bb$, where $+$ denotes the usual sum of two matrices.
\end{exercise}

\begin{sol}
    If $P, Q \in \bb$, then $M(P + Q) = MP + MQ = PM + QM = (P + Q)M$, so $P + Q \in \bb$.
\end{sol}

\begin{exercise}[0.1.3]
    Prove that if $P, Q \in \bb$, then $P\cdot Q \in \bb$, where $\cdot$ denotes the usual product of two matrices.
\end{exercise}

\begin{sol}
    If $P, Q \in \bb$, then $M(PQ) = (MP)Q = (PM)Q = P(MQ) = P(QM) = (PQ)M$, so $PQ \in \bb$.
\end{sol}

\begin{exercise}[0.1.4]
    Find conditions on $p, q, r, s$ which determine precisely when
    \[
        \begin{pmatrix}
            p & q \\
            r & s
        \end{pmatrix} \in \bb.
    \]
\end{exercise}

\begin{sol}
    Let $X$ be the matrix above.
	Computing $MX = XM$ gives the following system of equations:
    \[
        \begin{cases}
            p + r = p \\
            q + s = p + q \\
            r = r \\
            s = r + s
        \end{cases}
    \]
    Simplifying results in $r = 0$ and $p = s$.
	Then
    \[
        \bb = 
        \set*{\begin{pmatrix}
            p & q \\
            0 & p
        \end{pmatrix}~\middle|~p, q \in \r}. \qh
    \]
\end{sol}

\begin{exercise}[0.1.5]
    Determine whether the following functions $f$ are well-defined:
    \begin{subexercises}
        \item $f : \q \to \z$ defined by $f(a/b) = a$.
        \item $f : \q \to \q$ defined by $f(a/b) = a^2/b^2$.
    \end{subexercises}
\end{exercise}

\begin{solenum}
    \item No: $1/2 = 2/4$, but $f(1/2) = 1 \neq 2 = f(2/4)$.
    \item Yes: $a/b = c/d$ implies $a^2/b^2 = c^2/d^2$, so $f(a/b) = f(c/d)$.
\end{solenum}

\begin{exercise}[0.1.6]
    Determine whether the function $f: \r^+ \to \z$ defined by mapping a real number $r$ to the first digit to the right of the decimal point in a decimal expansion of $r$ is well-defined.
\end{exercise}

\begin{sol}
    No.
	Note that $1=1.000\ldots = 0.999\ldots$, but $f(1.000\ldots) = 0$ and $f(0.999\ldots) = 9$.
\end{sol}

\begin{exercise}[0.1.7]
    Let $f: A \to B$ be a surjective map of sets.
	Prove that the relation
    \[
    	a \sim b \iff f(a) = f(b)
    \]
    is an equivalence relation whose equivalence classes are the fibers of $f$.
\end{exercise}

\begin{sol}
    That $\sim$ is an equivalence relation follows from $=$ being an equivalence relation.

    The equivalence class of $a \in A$ is $\set{x \in A \mid x \sim a}$.
	By definition, this is $\set{x \in A \mid f(x) = f(a)}$, which is precisely the fiber of $f$ over $f(a)$.
\end{sol}